\subsection{Short composite hom-sets and vanishing products}
\begin{proposition}[Proved]\label{A7:short-composite-gates}
Let $C$ be a finite acyclic category with support $0<1<2<3<4$,
and let $K$ be any field.
Write $A=C(0,1)$, $B=C(1,2)$, $D=C(2,3)$, $E=C(3,4)$, and
\[
 u(a,b)=b\circ a,\qquad v(b,c)=c\circ b,\qquad
 w(c,d)=d\circ c.
\]
Every product $H^2(BC;K)\otimes H^2(BC;K)\to H^4(BC;K)$ is zero
if any one of the following conditions holds:
\begin{enumerate}
\item for each $b\in B$, $u(a,b)$ is independent of $a\in A$;
\item for each $c\in D$, $w(c,d)$ is independent of $d\in E$;
\item there is $c\in D$ such that $v(b,c)$ is independent of $b\in B$;
\item there is $b\in B$ such that $v(b,c)$ is independent of $c\in D$.
\end{enumerate}
In particular, nonzero degree-two products require
$|C(0,2)|,|C(1,3)|,|C(2,4)|\ge2$.
\end{proposition}
\begin{proof}
Suppose (1) holds. Fix $a_*\in A$ and two cocycles $\alpha,\beta$.
For $a\ne a_*$ put
\[
 \ell_{ab}=\alpha(a,b)-\alpha(a_*,b),\qquad
 F_a(s)=\alpha(a,s)-\alpha(a_*,s),\quad s\in C(1,3).
\]
Subtracting the cocycle equations on $(a,b,c)$ and $(a_*,b,c)$
gives $\ell_{ab}=F_a(v(b,c))$ for every $b,c$.
For any four-cycle $z$, the face omitting object $0$ gives a unique
expression
\[
 z=\sum_{a\ne a_*}(a-a_*)\otimes
       \sum_{b,c,d}t^a_{bcd}\,b\otimes c\otimes d.
\]
Because the five faces of a four-simplex lie on different object
subsets, they vanish separately. The faces omitting $2$, $3$, and
$4$ give, for each $a\ne a_*$,
\[
 \sum_{b,c}t^a_{bcd}[v(b,c)]=0\quad(d\in E),\qquad
 \sum_{c,d}t^a_{bcd}[w(c,d)]=0\quad(b\in B),\qquad
 \sum_dt^a_{bcd}=0\quad(b,c).
\]
Here brackets denote the basis of the free $K$-vector space on the
indicated hom-set. The cocycle equation on $(b,c,d)$ gives
\[
 \beta(c,d)=\beta(v(b,c),d)-\beta(b,w(c,d))+\beta(b,c).
\]
Consequently the cup evaluation is
\[
 \langle\alpha\smile\beta,z\rangle
   =\sum_{a,b,c,d}t^a_{bcd}\ell_{ab}\beta(c,d)=0.
\]
Indeed, after the substitution, the first sum vanishes by the first
boundary equation weighted by $F_a(s)\beta(s,d)$; the second by
the second equation weighted by $\ell_{ab}\beta(b,r)$; and the
third by the last equation. Over a field, evaluation on all cycles
detects cohomology, proving (1). Reversing the objects proves (2);
the reversal interchanges the factors, whose degree is two.

For (3), fix such a $c$ and write $q=v(b,c)$, independent of $b$.
Define a one-cochain on the full subcategory on $012$ by
\[
 \lambda(a)=\alpha(a,q),\qquad
 \lambda(b)=\alpha(b,c),\qquad
 \lambda(p)=\alpha(p,c)\quad(p\in C(0,2)).
\]
The cocycle equation on $(a,b,c)$ says exactly
$\alpha(a,b)=\lambda(b)-\lambda(u(a,b))+\lambda(a)$.
Thus $\alpha$ restricts to a coboundary on $012$. Extending
$\lambda$ by zero to all other edges and subtracting its coboundary
gives a representative vanishing on every $012$ triangle. Its
Alexander--Whitney product with $\beta$ is zero on every
four-simplex. Reversal proves (4). The three stated cardinality
conditions follow from (1), (3), and (2), respectively.
\end{proof}

\begin{proposition}[Proved]\label{A7:long-composite-gates}
For a finite acyclic category with support $0<1<2<3<4$, nonzero
degree-two products over a field require that at least two of
$C(0,3),C(0,4),C(1,4)$ have size at least two.
\end{proposition}
\begin{proof}
If $C(0,4)$ and $C(1,4)$ are singletons, choose any arrow
$r:2\to4$. The unique arrows $0\to4$, $1\to4$, together with
$r$, are the components of a natural transformation from the
inclusion of the full subcategory on $012$ to the constant functor
at $4$. Naturality follows from the two singleton assumptions.
Thus the inclusion is null-homotopic on classifying spaces, so
every two-cocycle restricts to a coboundary on $012$.
Extending a primitive by zero shows that every degree-two product
is zero by the Alexander--Whitney formula. Opposition treats
singleton $C(0,3)$ and $C(0,4)$.

It remains that $C(0,3)=\{t\}$ and $C(1,4)=\{r\}$.
Fix arrows $c:2\to3$ and $d:3\to4$. For any two-cocycle $\alpha$,
the equation on $(a,s,d)$, where $a:0\to1$ and $s:1\to3$,
is
\[
 \alpha(a,s)=\alpha(s,d)-\alpha(t,d)+\alpha(a,r).
\]
Substitute $s=c\circ b$ in the cocycle equation on $(a,b,c)$.
This gives
\[
 \alpha(a,b)=\alpha(b,c)+\alpha(c\circ b,d)-\alpha(t,d)
             -\alpha(b\circ a,c)+\alpha(a,r).
\]
The right side is $\delta\lambda(a,b)$ for the one-cochain on
$012$ defined by
\[
 \lambda(a)=\alpha(a,r),\quad
 \lambda(b)=\alpha(b,c)+\alpha(c\circ b,d)-\alpha(t,d),\quad
 \lambda(p)=\alpha(p,c)\quad(p:0\to2).
\]
Its extension again kills the restriction to $012$ and therefore
every degree-two product. These are all pairs among the three
indicated hom-sets.
\end{proof}

\begin{proposition}[Proved: collapse of the longest hom-set]
\label{A7:extreme-arrow-collapse}
Let $C$ be a finite acyclic category on $0<1<2<3<4$ with total-chain
support, and let $C^{\flat}$ be obtained by replacing $C(0,4)$ by
one arrow and retaining every other arrow. The canonical quotient
functor $q:C\to C^{\flat}$ has
\[
 q^*:H^2(BC^{\flat};F)\twoheadrightarrow H^2(BC;F),
 \qquad
 q^*:H^j(BC^{\flat};F)\xrightarrow{\cong}H^j(BC;F)
 \quad(j\geq3)
\]
for every field $F$, and preserves cup products. Consequently, if
$H^2(BC;F)\otimes H^2(BC;F)\to H^4(BC;F)$ is nonzero, then
$|C(0,3)|\geq2$ and $|C(1,4)|\geq2$.
\end{proposition}
\begin{proof}
Keep all original compositions whose target hom-set is not $C(0,4)$,
and send every composition into the new singleton hom-set to its
unique arrow. This is an associative category: every associativity
equation with target $C(0,4)$ becomes an equality of that unique
arrow, and every other equation is unchanged. The map $q$ is the
identity on objects and all other arrows, and sends all $0\to4$
arrows to the new arrow. It is a functor and reduces the arrow count
by $|C(0,4)|-1$.

An arrow $0\to4$ cannot occur as a factor of a nondegenerate
simplex of degree at least two: there is no nonidentity arrow into
$0$ or out of $4$. Thus $q$ identifies the normalized chain bases
in every degree at least two. The boundary $d_2$ can change,
because a composite of two arrows may lie in $C(0,4)$, but
$d_j$ is identical for $j\geq3$. On cochains, therefore, $q^*$
is the identity in degrees at least two and intertwines identical
coboundaries there. The degree-two cocycles coincide, while
\[
 \operatorname{im}\bigl(\delta^1_{C^{\flat}}\bigr)
 \subseteq\operatorname{im}\bigl(\delta^1_C\bigr),
\]
where the inclusion follows from the cochain-map identity for
$q^*$. Taking quotients gives the asserted surjection in degree
two and isomorphisms in degrees at least three. Functoriality
of the Alexander--Whitney product gives compatibility with cups.
No assertion that $q$ is a homotopy equivalence or a degree-two
homology equivalence is used.

If either $C(0,3)$ or $C(1,4)$ is a singleton, then $C^{\flat}$
has that singleton hom-set and the singleton $C^{\flat}(0,4)$.
Proposition~\ref{A7:long-composite-gates} makes all its degree-two
products zero. Lift any two degree-two classes of $BC$ along the
proved surjection. Their product is the pullback of zero, hence
zero. This proves the final assertion.
\end{proof}
